%% kapa: exam=1 round=22 year=2011 period=2 form=A subject=appraisal_law subject_ko=부동산관계법규 seq=28 number=68
%% source: Q-Net 감정평가사 기출문제 (cst00309) artlSeq=5036790 file=extracted/1008/5036790/2011년 제22회 감정평가사 1차시험 기출문제/2교시.hwp
%% conversion: kordoc 4.24.51 -> tools/gen1.py (auto); review=proofread-gpt-astra; flags=none
\begin{question}{68}
국유재산법령상 국유재산에 관한 설명으로 옳지 않은 것은?
\begin{choices}
  \item 국유재산에 관한 사무에 종사하는 직원이 소관청의 허가를 받지 아니하고 자신이 처리하는 국유재산을 취득한 경우 이는 무효로 한다.
  \item 국유재산의 관리ㆍ처분에 관한 소관 중앙관서의 장이 없는 경우에는 총괄청이 이를 지정한다.
  \item 등기가 필요한 국유재산의 경우 그 권리자의 명의는 총괄청의 명칭으로 한다.
  \item 행정재산은 시효취득의 대상이 되지 아니한다.
  \item 기부채납을 조건으로 하는 경우 사인이라도 국유하천에 교량 등 영구시설물을 축조할 수 있다.
\end{choices}
\end{question}
\answer{3}
